\documentclass[11pt,a4paper,sans]{moderncv}
\moderncvstyle{banking}
\moderncvcolor{blue}

% Force first/last name and section headings to render in moderncv blue (color1).
\renewcommand*{\firstnamestyle}[1]{{\fontsize{34}{36}\bfseries\upshape\color{color1}#1}}
\renewcommand*{\lastnamestyle}[1]{{\fontsize{34}{36}\bfseries\upshape\color{color1}#1}}
\renewcommand*{\sectionstyle}[1]{{\sectionfont\color{color1}#1}}

\usepackage{hyperref}
\hypersetup{
    colorlinks=true,
    linkcolor=blue,
    filecolor=magenta,
    urlcolor=blue,
    pdftitle={Walter Hernandez - CV},
    pdfpagemode=FullScreen,
}
\usepackage[scale=0.77]{geometry}
\usepackage{needspace}

\name{Walter}{Hernandez}
\address{Argentina}{}{}
\phone[mobile]{(+54) 11 2382-6514}
\email{walskp@gmail.com}
\extrainfo{\href{https://www.linkedin.com/in/wal-hernandez-dev}{LinkedIn}, \href{https://github.com/Wal-Hernandez}{GitHub}}

\begin{document}
\makecvtitle

\vspace{-4pt}

%------------------------------------------------------------------------------
\section{Profile}
%------------------------------------------------------------------------------
\vspace{1pt}
Software Engineer with hands-on experience across backend services, API design and AWS
cloud infrastructure. At La Naci\'{o}n I contributed to the migration from a legacy .NET 2.2 newsletter system to
a Node.js + AWS event-driven architecture, building REST APIs, data models and serverless backend microservices
with Lambda, EventBridge and SQS. I then joined the Paywall monorepo as a core developer, contributing
to business-critical features across the full stack. At Eppical I built backend integrations for
AI document-processing workflows using OpenAI APIs. I apply SOLID principles and design
patterns in production code and work across the full development lifecycle including CI/CD
pipelines, code reviews and technical documentation.

%------------------------------------------------------------------------------
\section{Professional Experience}
%------------------------------------------------------------------------------
\vspace{1pt}

\needspace{5\baselineskip}
\cventry{Oct 2024--Present}{Full Stack Developer SSR}{La Naci\'{o}n}{Argentina}{Hybrid/Remote}{\vspace{3pt}
\begin{itemize}
\item Contributed to the migration from a legacy .NET 2.2 newsletter system to Node.js + AWS, building
REST APIs, data models and serverless backend services integrated via EventBridge and SQS.
Configured Dead Letter Queues for fault tolerance and monitored service health with CloudWatch.
\item Apply SOLID principles and design patterns (Strategy, Factory, Repository, Dependency Injection)
in production backend code and API architecture.
\item Contributed to the Paywall monorepo as a core developer, delivering business-critical
features end to end and participating in technical refinement for complex backend tasks.
\item Work directly with Product Owners and business teams to clarify requirements, then split
large stories into concrete engineering tasks with Architecture, Database and Frontend teams.
\item Perform code reviews via Azure DevOps and maintain technical documentation in collaboration with
cross-functional teams.
\end{itemize}}

\vspace{3pt}
\cventry{Jul 2023--Oct 2024}{Full Stack Software Developer}{Eppical}{Argentina}{Remote}{\vspace{3pt}
\begin{itemize}
\item Built backend integrations for AI document-processing workflows, routing documents through
OpenAI APIs and processing analysis results in Node.js/.NET services consumed by Angular frontends.
\item Developed full-stack features across Angular frontends and Node.js/.NET backends, including UI
components, API endpoints and database interactions.
\item Refactored code to improve maintainability and participated in technical discussions to evaluate
implementation trade-offs with the team and Product Owners.
\end{itemize}}

\vspace{3pt}
\cventry{Feb 2021--Jul 2023}{Full Stack Developer}{SYNAgro}{Argentina}{On-site/Remote}{\vspace{3pt}
\begin{itemize}
\item Developed internal business applications for the agribusiness sector across the full stack, building
React and Ionic interfaces, implementing C\#/.NET backend services and writing SQL/database changes.
\item Built responsive web and mobile interfaces with React, TypeScript and Ionic, fixed production bugs
and implemented customer-facing features.
\end{itemize}}

%------------------------------------------------------------------------------
\section{Core Competencies}
%------------------------------------------------------------------------------
\vspace{1pt}
\begin{itemize}
\item \textbf{Backend \& APIs:} Node.js, TypeScript, C\#, .NET, REST API design, event-driven
architecture, backend microservices, SOLID principles, design patterns (Strategy, Factory,
Repository, DI), legacy modernization, API integrations.
\item \textbf{Cloud \& Infrastructure:} AWS (Lambda, API Gateway, EventBridge, SQS, CloudWatch),
Azure, Azure DevOps, Git Flow, CI/CD pipelines.
\item \textbf{Databases:} PostgreSQL, SQL Server, MongoDB.
\item \textbf{Frontend:} React, Angular, Ionic, Vue.js.
\item \textbf{AI \& Automation:} OpenAI APIs, document-analysis pipelines, AI-assisted development.
\item \textbf{Engineering Practices:} Cross-functional collaboration, code reviews, technical documentation,
Product Owner and business-team collaboration, Agile.
\end{itemize}

%------------------------------------------------------------------------------
\section{Education}
%------------------------------------------------------------------------------
\vspace{1pt}

\needspace{5\baselineskip}
\cventry{2024--Present}{Seguridad Inform\'{a}tica (In Progress)}{Universidad Siglo 21}{Argentina (Remote)}{}{\vspace{1pt}
Focus: infrastructure, security, and modern technology foundations.
}

\vspace{3pt}
\cventry{2021--2024}{Information Systems Engineering (Partial, Reoriented)}{National Technological University (UTN)}{Argentina}{}{\vspace{1pt}
Relevant coursework: Mathematics, Logic, Data Structures, Algorithms, Software Engineering.
}

\vspace{3pt}
\cventry{2021--2022}{Full Stack Web Development Bootcamp}{Henry}{Remote}{}{\vspace{1pt}
1500+ hours of intensive training, daily pair programming, individual and group projects.
}

%------------------------------------------------------------------------------
\section{Languages}
%------------------------------------------------------------------------------
\vspace{1pt}
\begin{itemize}
\item \textbf{Spanish:} Native
\item \textbf{English:} C1 Advanced
\end{itemize}

\end{document}
